\documentclass[10pt,a4paper]{amsart}
\usepackage[margin=19mm]{geometry}
\usepackage{amsmath,amssymb,mathtools}
\usepackage[scr=rsfs,cal=euler]{mathalfa}
\usepackage{xfrac}
\usepackage[unicode,hidelinks,bookmarks=false]{hyperref}
\newcommand{\ie}{i.e.\ }
\newcommand{\cf}{cf.\ }
\newcommand{\resp}{respectively\ }
\newcommand{\Subsection}{\subsection}
\providecommand{\nobreakdash}{\nobreak\hbox{-}}
\setlength{\emergencystretch}{2em}
\multlinegap0pt
\usepackage{mathtools}
\usepackage{bm}
\usepackage{tikz}
\usepackage{float}
\usepackage{amssymb}
\usepackage{tcolorbox}
\newtheorem{remark}{Remark}
\newtheorem{lemma}{Lemma}
\newtheorem{prop}{Proposition}
\newcommand{\C}{{\mathbb C}}
\newcommand{\CJ}{{\mathcal J}}
\newcommand{\CK}{{\mathcal K}}
\newcommand{\fG}{\mathfrak{G}}
\newcommand{\fg}{\mathfrak{g}}
\newcommand{\fh}{\mathfrak{h}}
\newcommand{\fn}{\mathfrak{n}}
\newcommand{\tH}{{\tilde H}}
\newcommand{\tL}{{\tilde L}}
\title{Infinite Dimensional Topological-Holomorphic Symmetry in Three-Dimensions}
\author{Hank Chen, Joaquin Liniado}
\date{Source: arXiv:2507.01858v4 (CC BY 4.0)}
\begin{document}
\maketitle

\noindent\emph{Source and licence.} Infinite Dimensional Topological-Holomorphic Symmetry in Three-Dimensions. Version arXiv:2507.01858v4 (CC BY 4.0), distributed by arXiv under the Creative Commons Attribution 4.0 International licence, http://creativecommons.org/licenses/by/4.0/, as recorded in the arXiv licence metadata for this exact version and shown on the abstract page https://arxiv.org/abs/2507.01858v4. Full licence notice: \texttt{licenses/arxiv-2507.01858-CC-BY-4.0.txt}.\par\smallskip

\noindent\textit{Editorial scope.} Counted unit: the article's prop prop:rectheorem (source0.tex lines 1760--1788). The article's complete original proof of it is delivered with it (source0.tex lines 1909--2075, one \texttt{proof} environment of 167 lines, which the article places away from the statement). No mathematics is written, repaired or re-derived: the statement and the proof are the article's own lines, in the article's own notation, and no retained line is substituted or re-flowed.  Retained uncounted as the source-local context the proof needs: lines 1305--1309; lines 1316--1318; lines 1380--1382; lines 1419--1422; lines 1710--1716; lines 1837--1840.  Nothing was omitted from any retained span, so the package carries no omission record.  No retained line contains a citation, so the wrapper carries no bibliography and no bibliography entry.  No retained line contains a figure environment, an image-inclusion command or drawing code, so no image is delivered and the package declares no asset dependency.  Editorial formatting only, in addition: an AMS \texttt{amsart} preamble that restates the article's own declarations and macros used by the retained lines, namely the environments \texttt{remark}, \texttt{lemma}, \texttt{prop} on separate counters, together with the standard packages those lines need.  The retained environments print in this document as Remark 1, Lemma 1, Proposition 1 in order of appearance, whereas the article prints the same items with the numbering of its own full document.  The complete native source of the article as distributed in the arXiv e-print for this exact version is bundled unchanged as \texttt{sources/180908/source0.tex}.
\begin{equation}
    [t^a, t^b]^\vee = \tilde f^{ab}_c\, t^c \in V^\vee, \qquad
    \mu_2^\vee(t^a,s^b) = (\tilde\mu_2)^{ab}_c\, s^c \in \mathfrak{n}^\vee\,,
    \label{ec:aff2liealgbrack}
\end{equation}

\begin{remark}\label{dualdegree}
    Given a Lie 2-algebra $\fG = \fh\xrightarrow{\mu_1}\fg$, the dual Lie 2-algebra $\fG^\vee[1] = \fg^\vee\xrightarrow{\mu_1^\vee}\fh^\vee$ is given by the underlying dual complex together with a shift in degree \textit{down} by $1$, denoted by "$[1]$", such that $\fg^\vee$ has degree $(-1)$ and $\fh^\vee$ has degree $0$. This is to ensure that $\mu_1^\vee$ remains a cohomological differential which increases the degree by $1$, and hence $\fG^\vee[1]$ can still be treated as a Lie 2-algebra.
\end{remark} 

\begin{equation}\label{degconvention}
     \text{deg}\left(V^\vee \otimes \CK_u^0\right) =0,\qquad \text{deg}\left(\fn^\vee \otimes \CK_u^1\right)=1\,.
\end{equation}

\begin{equation}
\label{ec:kilform}
    \langle t^a,s^b\rangle^\vee = \big(\kappa^{-1}\big)^{ab} \,.
\end{equation}

\begin{lemma}
\label{ec:equivmutualloc}
    The raviolo fields $A(z)$ and $B(w)$ are mutually local if and only if there exist raviolo fields $C_n(w)$ such that 
    \begin{equation}
        [A(z),B(w)]=\sum_{n=0}^N \frac{1}{n!}\partial_w^n\Delta(z-w)C_n(w)
    \end{equation}
\end{lemma}

\begin{equation}
\label{ec:vacummodule}
   \mathcal{V} = U\widehat{\mathfrak{G}}^\vee_{\mathfrak{w}} \otimes_{U\widehat{\mathfrak{G}}^\vee_{\geq 0}} \mathbb{C}_w\,, 
\end{equation}

\begin{prop}
\label{prop:rectheorem}
Let $\mathcal V = \bigoplus \mathcal V^r$ be a $\mathbb Z$--graded vector space, 
$|0\rangle$ a non--zero vector, $\partial$ a degree $0$ endomorphism of $\mathcal V$. 
Further, let $\{a^i\}$ be a countable ordered set of vectors in $\mathcal V$, 
with $a^i \in \mathcal V^{r_i}$, and suppose we are given homogeneous fields
\begin{equation}
    A^i(z) = \sum_{m\geq 0} z^{m} A^{i,+}_m +  \Omega^m_z A^{i,-}_m
\end{equation}
of degree $r_i$, such that the following hold:
\begin{enumerate}
    \item $A^{i,+}_{0}|0\rangle = a^i$ and $A^{i,-}_m |0\rangle = 0$ for all $i$ and $m \geq 0$.
    \item $\partial |0\rangle = 0$ and $[\partial, A^i(z)] = \partial_z A^i(z)$ for all $i$.
    \item All fields $A^i(z)$ are mutually local.
    \item $\mathcal V$ is spanned by the vectors
    \begin{equation}
        A^{i_1}_{j_1} \cdots A^{i_l}_{j_l} |0\rangle, \qquad j_k \geq 0\,.
    \end{equation}
\end{enumerate}
Then the assignment
\begin{multline}
   Y\!\left(A^{i_1}_{j_1}\cdots A^{i_l}_{j_l}|0\rangle, z\right) \\
= \frac{1}{(j_1-1)! \cdots (j_l-1)!} 
:\partial_z^{-j_1-1} A^{i_1}(z) \cdots \partial_z^{-j_l-1} A^{i_l}(z): 
\end{multline}
determines a well-defined raviolo vertex algebra structure on $\mathcal V$. 
Moreover, this is the unique raviolo vertex algebra structure on $\mathcal V$ 
satisfying conditions (1)--(4) and such that $Y(a^i, z) = A^i(z)$.
\end{prop}

\begin{proof}
Let us take the vector space $\mathcal{V}$ defined in \eqref{ec:vacummodule} and consider, for $1\leq a\leq \operatorname{dim}\mathfrak{n}^\vee=\operatorname{dim}V^\vee$, the raviolo field operators
\begin{equation}
    \mathcal{J}^a(z) = \sum_{n\geq 0} z^n \tL_n^a  + \Omega_z^n \tH^a_n
\end{equation}
which, with the degree conventions in \textit{Remark \ref{dualdegree}}, are homogeneous of degree $1$. We identify the unit $ |0\rangle=1\otimes 1\in\mathcal{V}$ as the vacuum. We will now prove points (1) through (4) of \textbf{Proposition} \ref{prop:rectheorem}. 



By definition of $\mathcal{V}$, we have that
\begin{equation}
    \tH_n^a |0\rangle=0 \,,\quad  \tL_0^a|0\rangle = \ell_0^a\,,\label{vacmodulepolynomials}
\end{equation}
which proves point (1). Next, for the translation operator, we 
begin by defining an additional generator $D$ on $\widehat{\fG}^\vee_\mathfrak{w}\cong (\fG^\vee\otimes\mathcal{K}^\bullet_u)\oplus \C[-1]$ from the action of the derivative $\partial_u$ on the current modes:
\begin{equation*}
    [D,X\otimes k(u)] = X\otimes\partial_u k(u), \qquad [D,\mathfrak{w}]=0\,,
\end{equation*}
where $X\in\fG^\vee,~ k(u)\in\mathcal{K}_u^\bullet$ and $\mathfrak{w}\in\C[-1]$ is the degree 1 generator.  Under the canonical $\widehat{\fG}^\vee_\mathfrak{w}$-representation  on $\mathcal{V}$, this gives rise to an endomorphism $\partial\in\operatorname{End}\mathcal{V}$ which, by construction, satisfies
\begin{equation}
    [\partial,x]=\partial_z x\,,\qquad \partial|0\rangle = 0\label{translate}
\end{equation}
for all $x\in \widehat{\fG}^\vee_\frak{w}$. More explicitly,  these read as the following mode relations
\begin{equation}
\label{ec:actionpartial1}
    [\partial ,\tL_n^a] = (n+1)\tL_{n+1}^a\,,\quad [\partial, \tH_n^a] = -n \tH_{n-1}^a\,.
\end{equation}
This then allows us to compute
\begin{equation}
\label{ec:actionpartial2}
\begin{split}
    [\partial,\mathcal{J}^a(z)]
    &=\sum_{n\geq 0}z^n[\partial,\tL_n^a]+ \Omega_z^n[\partial,\tH_n^a] \\
&\stackrel{\eqref{ec:actionpartial1}}{=} \sum_{n\geq 0 }(n+1)z^n \tL_{n+1}^a - n\Omega^{n}_z\tH_{n-1}^a=\partial_z \mathcal{J}^a(z)\,,
\end{split} 
\end{equation}
identifying $\partial$ as the desired translation operator, thereby proving point (2). 

\medskip

To prove point (3), mutual locality, we will directly compute the commutator $[\mathcal{J}^a(z),\mathcal{J}^b(w)]$ between two raviolo fields and show that it is of the form in \textbf{Lemma \ref{ec:equivmutualloc}}. That is, we will show that 
\begin{equation}
\label{ec:point3aim}
    [\mathcal{J}^a(z),\mathcal{J}^b(w)]= (\bm{\kappa}^{-1})^{ab}\partial_w\Delta(z-w) + \mathbf{\tilde f}^{ab}_c\Delta(z-w)\mathcal{J}^c(w)\,.
\end{equation}
This computation involves a few grading-related technicalities, so we will present it in full detail.


To begin with, although we have been using the same Lie algebra index for both $\tL$  and $\tH$, in fact these fields take values in two distinct vector spaces. It will therefore be important to distinguish them in what follows. We thus introduce a collective index $a=(A,\alpha)$ so that
\begin{equation}
    \mathcal{J}^{a}(z)=\sum_{n\geq 0}z^n\tL^A_n + \Omega_z^n\tH^\alpha_n \,.
\end{equation}
Furthermore, recall from \eqref{degconvention} that $\tilde L_n^A$ has degree $1$ while $\tilde H_n^\alpha$ has degree $0$, and that $z^n$ has degree zero and $\Omega_z^n$ has degree $1$. As such, the raviolo field $\mathcal{J}^a(z)$ has homogeneous degree $1$ in $\mathcal{K}^\bullet_z \otimes \mathfrak{G}^\vee$. The commutator can thus be computed as
\begin{multline}
\label{ec:finalcom}
    [\CJ^a(a),\CJ^b(w)]=\sum_{n,m\geq 0} z^nw^m[\tL_n^A,\tL_m^B]-z^n\Omega^m_w[\tL_n^A,\tH_m^\beta]\\
    +\Omega_z^nw^m[\tH_n^\alpha,\tL_m^B]+\Omega_z^n\Omega_w^m[\tH_n^\alpha,H_m^\beta]\,,
\end{multline}
where a minus sign appeared in the second term from commuting $\tilde L_n^A$ past $\Omega_w^m$, both of which have odd degree $1$.\footnote{We thank Niklas Garner for pointing this out.} 


At this point we must be careful when replacing the mixed commutators in \eqref{ec:finalcom} by their explicit expressions, which we recall here for convenience:
\begin{equation}
\label{ec:higherpossbrack4}
[\tH_n^\alpha,\tL_m^B]
=
\begin{cases}
(\tilde \mu_2)^{\alpha B}_C\,\tL_{m-n}^C
+ n(\kappa^{-1})^{\alpha B}\,\delta_{n-1,m}
& \text{if } m \ge n-1 \,,\\[0.3em]
0 & \text{otherwise} \,.
\end{cases}
\end{equation}
Recall that $(\tilde\mu_2)^{\alpha B}{}_C$ and $(\kappa^{-1})^{\alpha B}$ were defined in
\eqref{ec:aff2liealgbrack} and \eqref{ec:kilform}, respectively, as the structure constants of the dual maps
\begin{equation}
  \mu_2^\vee : V^\vee \times \fn^\vee \to \fn^\vee,
  \qquad
  \langle\,\cdot\,,\,\cdot\,\rangle^\vee : V^\vee \times \fn^\vee \to \C .
\end{equation}
Hence, in order to write the commutator
$[\tL_n^A,\tH_m^\beta]$ appearing in \eqref{ec:finalcom}, we introduce
the map
\begin{equation}
\nu_2^\vee : \fn^\vee \times V^\vee \to \fn^\vee \,,
\qquad
\nu_2^\vee(s^A,t^\beta)
= -\,\mu_2^\vee(t^\beta,s^A)\,,
\end{equation}
so that, in components,
\[
(\tilde \nu_2)^{A\beta}_C
= -(\tilde \mu_2)^{\beta A}_C \,.
\]
In complete analogy, we define
\begin{equation}
\eta^{-1} : \fn^\vee \times V^\vee \to \mathbb{C} \,,
\qquad
\eta^{-1}(s^A,t^\beta)
= \kappa^{-1}(t^\beta,s^A)\,,
\end{equation}
such that
\[
(\eta^{-1})^{A\beta}
= (\kappa^{-1})^{\beta A} \,.
\]
These definitions are chosen precisely so as to ensure the skew-symmetry of the bracket $[\tL_n^A,\tH_m^\beta]=-[\tH_m^\beta,\tL_n^A]$. With these conventions in place, we may rewrite \eqref{ec:finalcom} as
\begin{multline}
\label{ec:expparacomp}
        [\CJ^a(z),\CJ^b(w)]=\sum_{n,m\geq 0}-z^n\Omega_w^m\left[(\tilde{\nu}_2)^{A\beta}_C\tL_{n-m}^C + m(\eta^{-1})^{A\beta }\delta_{m-1,n}\right]\\
        +\Omega_z^nw^m\left[(\tilde{\mu}_2)^{\alpha B}_C\tL_{m-n}^C + n(\kappa^{-1})^{\alpha B}\delta_{n-1,m}\right] + \Omega_z^n\Omega_w^m \tilde{f}^{\alpha\beta}_\gamma \tH_{n+m}^\gamma\,.
\end{multline}


\medskip

Turning now to the right-hand side of equation \eqref{ec:point3aim}, we compute 
\begin{equation}
    \Delta(z-w)\mathcal{J}^c(w)=\sum_{n,m\geq 0}w^{n+m}\Omega_z^n\tL_m^C + \Omega_z^n\Omega_{w}^{m-n}\tH_m^\gamma -z^n\Omega_w^{n-m}\tL_m^C
\end{equation}
which is an element of 
\begin{equation}
\label{ec:dirsumdec1}
    (\mathcal{K}_w^0\otimes \mathcal{K}_z^1\otimes \fn^\vee)\oplus(\mathcal{K}_w^1\otimes \mathcal{K}_z^1\otimes V^\vee)\oplus(\mathcal{K}_w^1\otimes \mathcal{K}_z^0\otimes \fn^\vee)\,.
\end{equation}
We define the linear map $\mathbf{\tilde f}^{ab}_c$ acting on this space as the endomorphism:
\begin{equation}
    \mathbf{\tilde f}^{ab}_c=\iota_1 \circ (\tilde{\mu}_2)^{\alpha B}_C \circ p_1 + \iota_2 \circ \tilde f^{\alpha \beta}_\gamma\circ p_2 + \iota_3 \circ (\tilde{\nu}_2)^{A\beta}_C \circ p_3
\end{equation}
where $\iota_i,p_i$ with $i=1,2,3$ are inclusions and projections, respectively, of the direct sum decomposition \eqref{ec:dirsumdec1}. By construction, we then have 
\begin{multline}
\label{ec:granf1}
    \mathbf{\tilde f}^{ab}_c\Delta(z-w)\mathcal{J}^c(w) = \sum_{n,m\geq 0}w^{n+m}\Omega_z^n(\tilde{\mu}_2)^{\alpha B}_C\tL_m^C + \Omega_z^n\Omega_{w}^{m-n}f^{\alpha \beta}_\gamma\tH_m^\gamma \\
    -z^n\Omega_w^{n-m}(\tilde{\nu}_2)^{A\beta}_C\tL_m^C\,.
\end{multline}
In complete analogy, we have that
\begin{equation}
    \partial_w\Delta(z-w)=\sum_{n\geq 0} nw^{n-1}\Omega_z^n + (n+1)z^n\Omega_w^{n+1}\,,
\end{equation}
which can be thought as an element of 
\begin{equation}
\label{ec:dirsumdec2}
    (\CK_w^0\otimes \CK_z^1)\oplus(\CK_w^1\otimes \CK_z^0)\,,
\end{equation}
so that we may define the contraction with $(\bm{\kappa}^{-1})^{ab}$ as a linear endomorphism of the above space given by 
\begin{equation}
    (\bm{\kappa}^{-1})^{ab}=\iota_1\circ (\kappa^{-1})^{\alpha B}\circ p_1 - \iota_2\circ (\eta^{-1})^{A\beta}\circ p_2\,,
\end{equation}
where again, $\iota_j,p_j$ with $j=1,2$ are inclusions and projections, respectively, of the direct sum decomposition \eqref{ec:dirsumdec2}. By definition, we thus obtain
\begin{equation}
\label{ec:grank1}
    (\bm{\kappa}^{-1})^{ab}\partial_w\Delta(z-w)=\sum_{n\geq 0} n(\kappa^{-1})^{\alpha B}w^{n-1}\Omega_z^n - (n+1)(\eta^{-1})^{A\beta}z^n\Omega_w^{n+1}\,.
\end{equation}
Putting together equations \eqref{ec:granf1} and \eqref{ec:grank1}, and comparing with \eqref{ec:expparacomp} we finally have 
\begin{equation}
    [\mathcal{J}^a(z),\mathcal{J}^b(w)]= (\bm{\kappa}^{-1})^{ab}\partial_w\Delta(z-w) + \mathbf{\tilde f}^{ab}_c\Delta(z-w)\mathcal{J}^c(w)\,
    \label{commutator}
\end{equation}
as desired.


\medskip

Finally, point (4) follows immediately by construction of $\mathcal{V}$, and thus by the reconstruction \textbf{Theorem \ref{prop:rectheorem}}, $\mathcal{V}$ defines a unique raviolo vertex algebra structure.


\end{proof}

\end{document}
